\clearpage
\appendix

\section{Additional Measurement Results from Commercial MEC Deployments}
\label{appnd:measurements}

This appendix presents additional measurement results from our commercial MEC deployment studies that complement the findings presented in \autoref{sec:motivation}.

\subsection{End-to-End Latency of Augmented Reality }
\label{appnd:ar}

\begin{figure}[h!]
    \centering
    \includegraphics[width=\columnwidth]{figures/measurement-contention-results/ran-contention/ar/e2e_latency_cdf.pdf}
    \caption{End-to-end latency for the augmented reality application without edge resource contention across MEC deployments in three cities. The dotted red line indicates the SLO.}
    \label{fig:motive-e2e-ar-ran}
\end{figure}

\autoref{fig:motive-e2e-ar-ran} shows that augmented reality (AR) exhibits a similar long-tail latency distribution as smart stadium during periods of low network activity. 
Since AR requires lower uplink throughput than smart stadium, only about \percent{5} of requests miss their SLO across all three cities. 
However, under high network activity (\dallas-Busy), wireless resource contention causes over \percent{98} of requests to violate their SLO.

\subsection{Effect of Compute Contention on End-to-End Latency}
\label{appnd:compute-contention}

We show the remaining results under compute resource contention for smart stadium in \nanjing and \seoul and augmented reality in all three cities.

\label{appnd:compute-contention-ss}
\begin{figure}[h!]
    \centering
    \includegraphics[width=\columnwidth]{figures/measurement-contention-results/compute-contention/video-transcoding/Nanjing/e2e_latency_cdf.pdf}
    \vspace{-2em}
    \caption{End-to-end latency for smart stadium under different levels of compute resource contention in \nanjing. The dotted red line indicates the SLO.}
    \label{fig:motive-e2e-transcoding-compute-nanjing}
\end{figure}

\begin{figure}[h!]
    \centering
    \includegraphics[width=\columnwidth]{figures/measurement-contention-results/compute-contention/video-transcoding/Seoul/e2e_latency_cdf.pdf}
        \vspace{-2em}
 
    \caption{End-to-end latency for smart stadium under different levels of compute resource contention in \seoul. The dotted red line indicates the SLO.}
    \label{fig:motive-e2e-transcoding-compute-seoul}
\end{figure}


\mypara{End-to-end latency of smart stadium.}
\autoref{fig:motive-e2e-transcoding-compute-nanjing} and \autoref{fig:motive-e2e-transcoding-compute-seoul} further confirm the effect of CPU resource contention on the end-to-end latency of CPU-intensive applications like smart stadium.
As CPU load increases, more requests exceed their SLO requirements and suffer from long tail latency.

\begin{figure}[h!]
    \centering
    \includegraphics[width=\columnwidth]{figures/measurement-contention-results/compute-contention/ar/Dallas/e2e_latency_cdf.pdf}
        \vspace{-2em}

    \caption{End-to-end latency for augmented reality under different levels of compute resource contention in \dallas.}
    \label{fig:motive-e2e-ar-compute-dallas}
\end{figure}

\begin{figure}[h!]
    \centering
    \includegraphics[width=\columnwidth]{figures/measurement-contention-results/compute-contention/ar/Nanjing/e2e_latency_cdf.pdf}
    \vspace{-2em}

    \caption{End-to-end latency for augmented reality under different levels of compute resource contention in \nanjing.}
    \label{fig:motive-e2e-ar-compute-nanjing}
\end{figure}

\begin{figure}[h!]
    \centering
    \includegraphics[width=\columnwidth]{figures/measurement-contention-results/compute-contention/ar/Seoul/e2e_latency_cdf.pdf}
    \vspace{-2em}

    \caption{End-to-end latency for augmented reality under different levels of compute resource contention in \seoul.}
    \label{fig:motive-e2e-ar-compute-seoul}
\end{figure}


\mypara{End-to-end latency of augmented reality.}
To further validate the impact of GPU contention on the end-to-end latency of GPU-intensive applications such as augmented reality, we run experiments in \dallas, \nanjing, and \seoul under varying levels of GPU load. 
We implement a GPU stressor using CUDA to emulate different contention levels. 
As shown in \autoref{fig:motive-e2e-ar-compute-dallas}, \autoref{fig:motive-e2e-ar-compute-nanjing}, and \autoref{fig:motive-e2e-ar-compute-seoul}, higher stress levels lead to more requests exceeding their SLOs and experiencing longer tail latencies.


\subsection{Variability of Network Latency} 
\label{appnd:latency-variability}

\begin{figure*}[]
    \centering
    \begin{subfigure}[b]{0.48\textwidth}
        \centering
    \includegraphics[width=\columnwidth]{figures/measurement-decompose-appendix/Nanjing/total_latency_boxplot.pdf}
    \caption{\nanjing}
    \label{fig:motive-decompose-nanjing}
    \end{subfigure}
    \hfill
    \begin{subfigure}[b]{0.48\textwidth}
        \centering
    \includegraphics[width=\columnwidth]{figures/measurement-decompose-appendix/Seoul/total_latency_boxplot.pdf}
    \caption{\seoul}
    \label{fig:motive-decompose-seoul}
    \end{subfigure}
    \caption{Network latency variability for uplink and downlink transmissions across different data sizes in \nanjing and \seoul.}
    \label{fig:motive-decompose-combined}
\end{figure*}

\begin{minipage}{\columnwidth}
\autoref{fig:motive-decompose-combined} reveals consistent asymmetry in both \nanjing and \seoul. 
As request size grows, uplink latency shows pronounced variability, whereas downlink latency remains largely stable. 
This resembles our observation that uplink paths suffer from high variability, while downlink paths maintain relative stability.
\end{minipage}

\section{Deadline-aware Proactive Edge Resource Scheduling Algorithm}
\label{appnd:edge-scheduling-algo}

\begin{algorithm}[htbp]
\small
\DontPrintSemicolon
\caption{\modified{Deadline-aware proactive edge resource scheduling.}}
\label{alg:edge-proactive-sched}
\modified{\textbf{Notation:} $r$ denotes a request; $a$ denotes the application that serves $r$; $\texttt{type}(r) \in \{\textsc{cpu}, \textsc{gpu}\}$.\;}

\modified{\BlankLine
\If{$t^{edge}_{budget}(r) \le 0$}{
\texttt{drop}(r)\;
\Return\;
}

\BlankLine
$urgency(r) \leftarrow t^{edge}_{budget}(r) / SLO_a$\;

\BlankLine
\eIf{$\texttt{type}(r)=\textsc{cpu}$}{
    \If{$urgency(r) < 0.1$}{
        \If{\texttt{now} $- t_{last\_cpu\_alloc}(a) \ge \ms{100}$}{
        \texttt{assign\_one\_more\_core}(a)\;
        $t_{last\_cpu\_alloc}(a) \leftarrow \texttt{now}$\;
        }
    }
    \If{$\texttt{cpu\_util}(a) < 60\%$}{
    \texttt{reclaim\_one\_core}(a)\;
    }
}{
    $prio(r) \leftarrow \texttt{map\_urgency\_to\_prio}(urgency(r))$\;
    \texttt{set\_cuda\_stream\_priority}(r,\ prio(r))\;
}

\BlankLine
\texttt{process\_request}(r)\;
}
\end{algorithm}